\chapter{روش پیشنهادی و پیاده‌سازی}
\label{ch:method}

این فصل ساختار سامانه تحویل‌شده، زنجیره داده، آشکارساز وقفه، طراحی
آزمایش‌های مسیر مستقیم، و روش کاری شواهد‌محور را شرح می‌دهد.

\section{تلاش اولیه و بازنگری}
\label{sec:early-attempt}

نخستین پیاده‌سازی بر پایه \en{LLaMA-Omni2} بود: کدگذار منجمد
\en{Whisper}، لایه تصویرگر، مدل زبانی \en{Qwen2.5-0.5B}~\cite{qwen2025qwen25}
با آداپتر \en{LoRA}، و سرهای جداگانه بازشناسی و تولید. زیان کاهش یافت، اما
ممیزی درونی سه اشکال بنیادی را نشان داد: مدل زبانی در تابع پیش‌رو مشارکت
نداشت؛ هدف بازشناسی از درهم‌سازی‌های تصادفی در هر فرایند استفاده می‌کرد؛ و
هدف تولید، بازسازی همان اظهار ورودی بود نه پاسخ گفت‌وگویی. آن مصنوع یک مدل
گفتار‌به‌گفتار نبود. دو تصمیم از این ممیزی آمد: سرویس‌دهی به زنجیره آبشاری
با مؤلفه‌های واقعی منتقل شد، و برای مسیر مستقیم تنها مدلی پذیرفته شد که
آموزش‌دهنده رسمی داشته باشد.

\section{ساخت مجموعه داده}
\label{sec:data}

پیکره عمومی فارسی با برچسب وقفه در دسترس نبود. مجموعه داده از آرشیو داخلی
محتوای گفتاری چهار کانال مشتق شد. استفاده پژوهشی داخلی با تأیید استاد
راهنما ثبت شده است؛ این تأیید مجوز بازنشر صوت یا زیرنویس خام نیست و هیچ
داده محدودشده‌ای در مخزن کد نیست.

در طرح نخست، برچسب متنی با اجرای مجدد سرویس بازشناسی ساخته می‌شد. همان
سرویس پیش‌تر روی زیرنویس همین کانال‌ها تنظیم شده بود و یک حلقه دوری
می‌ساخت. پس از شناسایی این مشکل، برچسب متنی به زیرنویس اصلی منبع تغییر یافت
و بازشناسی تنها در نقش مؤلفه زمان اجرا و ابزار سنجش باقی ماند.

زنجیره چنین است: موجودی \N{775.887} ساعتی زیرنویس، انتخاب وزن‌دار برنامه،
بازسازی پنجره‌های پانزده‌دقیقه‌ای، تفکیک گوینده، برچسب شرایط نویز، استخراج
جفت «گفتار کاربر / پاسخ» بدون بازاستفاده، و ساخت استریوی ممیزی‌شده با کانال
صفر برای دستیار بازسازی‌شده و کانال یک برای کاربر. تفکیک
آموزش / اعتبارسنجی / آزمون بر شناسه نشست منبع است. برچسب‌های تعاملی
نامزد خودکارند و شمار ردیف تأییدشده انسانی صفر است.

\section{معماری سامانه آبشاری}
\label{sec:cascade}

سامانه تحویل‌شده یک زنجیره آبشاری با کنترل تمام‌دوطرفه بیرونی است
(شکل~\ref{fig:arch}). پیاده‌سازی بر \en{PyTorch}~\cite{paszke2019pytorch}
و \en{Transformers}~\cite{wolf2020transformers} استوار است.

\begin{figure}[ht]
\centering
\begin{latin}
\scalebox{0.88}{%
\begin{tikzpicture}[
  box/.style={draw,rounded corners=2pt,align=center,font=\scriptsize,
              minimum height=9mm,inner sep=3pt,text width=24mm},
  client/.style={box,fill=black!5},
  lbl/.style={font=\tiny,inner sep=1.5pt,fill=white},
  ar/.style={-{Latex[length=1.8mm]},thin},
  dashar/.style={ar,dashed}
]
\node[client] (mic) at (0,0) {continuous mic\\16 kHz PCM16};
\node[box] (feat) at (3.1,0)   {rolling features\\energy / $F_0$ / MFCC};
\node[box] (det)  at (6.2,0)   {GBDT barge-in\\classifier};
\node[box] (ctl)  at (9.3,0)   {playback\\controller};
\node[box] (asr)  at (3.1,-2.3) {NeMo Persian\\ASR};
\node[box] (llm)  at (6.2,-2.3) {Qwen3-4B\\prompt-v2};
\node[box] (tts)  at (9.3,-2.3) {Piper Persian\\TTS};
\node[client] (spk) at (12.4,-1.15) {browser audio\\source};
\draw[ar] (mic)  -- (feat);
\draw[ar] (feat) -- (det);
\draw[ar] (det)  -- (ctl);
\draw[ar] (asr)  -- (llm);
\draw[ar] (llm)  -- (tts);
\draw[ar] (mic.south) -- ++(0,-1.4) -- (asr.west)
      node[lbl,pos=0.62] {end of speech};
\draw[ar] (tts.east) -- (spk.south west)
      node[lbl,pos=0.5] {full PCM reply};
\draw[ar] (ctl.east) -- (spk.north west)
      node[lbl,pos=0.5] {stop source};
\draw[dashar] (ctl.north) -- ++(0,0.8) -| (asr.north)
      node[lbl,pos=0.3] {450 ms pre-roll becomes next turn};
\end{tikzpicture}}
\end{latin}
\caption{معماری سامانه تحویل‌شده. ردیف بالا مسیر کنترل تمام‌دوطرفه و ردیف
         پایین زنجیره تولید پاسخ است. پیکان نقطه‌چین، انتقال پیش‌بافر
         میکروفون به نوبت بعدی را نشان می‌دهد.}
\label{fig:arch}
\end{figure}

جریان نوبت کاربر به بازشناسی فارسی محلی می‌رود و از هنجارساز قطعی متن
می‌گذرد. پروتکل نهایی استفاده از کانال یک ممیزی‌شده را الزامی می‌کند؛ یک
پنل توسعه پیشین به‌اشتباه کانال صفر (صدای دستیار) را رونویسی می‌کرد و پیش
از هر ردیف اعتبارسنجی اصلاح شد.

پاسخ متنی با \en{Qwen3-4B-Instruct-2507}~\cite{qwen2025qwen3} به‌صورت محلی و
از درخت پرونده با درهم‌سازی تأییدشده تولید می‌شود. پرامپت سامانه قفل است.
نسخه \en{v2} سه قید دارد: پاسخ فارسی، کوتاه و گفتاری، و معطوف به آخرین
اظهار کاربر. اگر سهم نویسه فارسی از آستانه کمتر باشد یک تلاش دوباره انجام
می‌شود؛ اگر هر دو تلاش شکست بخورد، پاسخ جایگزین قاعده‌محور تولید می‌گردد.
شمار ردیف‌های جایگزین یکی از شروط پذیرش است.

پاسخ‌گوی تحویل‌شده همین مدل چهارمیلیاردپارامتری با پرامپت \en{v2} است. پنل
مکانیکی نه‌ردیفی و سنجه بازبازشناسی روی زنجیره پیش از ارتقا اجرا شده‌اند،
یعنی با \en{Qwen2.5-0.5B-Instruct} (بازنگری
\en{7ae557604adf67be50417f59c2c2f167def9a775}). آن پنل پس از ارتقا دوباره
اجرا نشد. آزمون معنایی \N{40} ردیفی روی \en{Qwen3-4B} است و داور آن
\en{Qwen3.8-27B-FP8} است (بازنگری
\en{017b9c7af6b5689d5dd426a76e0bc077eb5ca20a}).

متن پاسخ با \en{Piper} به موج تبدیل می‌شود. در همه ردیف‌های ارزیابی، شرط
پذیرش این است که موتور واقعاً \en{Piper} بوده باشد. صوت خروجی از نظر
متناهی‌بودن، طول کمینه، هنجارسازی دامنه و غیرساکت‌بودن با ریشه میانگین
مربعات بررسی می‌شود.

\section{تشخیص وقفه و بستر تمام‌دوطرفه}
\label{sec:bargein}

ویژگی‌ها روی جریان میکروفون غلتان محاسبه می‌شوند: فریم \N{20} میلی‌ثانیه،
پنجره \N{400} میلی‌ثانیه، \N{13} ضریب \en{MFCC} با دلتا، بسامد پایه
\en{YIN} در \N{60} تا \N{400} هرتز، انرژی لگاریتمی و نرخ گذر از صفر. برای
هر بعد، میانگین، انحراف معیار و بیشینه روی پنجره پیوند داده می‌شود.
طبقه‌بند تقویت گرادیانی \N{200} درخت، عمق \N{4}، نرخ یادگیری \N{0.08} و
کمینه \N{8} نمونه در برگ دارد. کلاس‌ها عبارتند از نبود رخداد، وقفه، بازخورد
کلامی و نویز. خط پایه، آستانه‌گذاری انرژی و نرخ گذر از صفر است.

کنترل‌گر پخش مالک تصمیم است. آشکارساز همیشه \N{450} میلی‌ثانیه آخر میکروفون
را نگه می‌دارد؛ هنگام پذیرش وقفه، این پیش‌بافر با صوت بعدی پیوند می‌خورد و
ورودی نوبت بعد را می‌سازد. بستر انتقال \en{WebSocket} است. آزمون‌های
بازگشتی، پیوستگی پیش‌بافر، تأیید توقف، لغو نوبت، اتصال مجدد، جداسازی نشست‌ها
و حالت نگه‌داری فقط‌سنجه را پوشش می‌دهند.

\مرزادعا{این آزمون‌ها پیوستگی را در سطح انتقال سرور اثبات می‌کنند. توقف
منبع صوت در مرورگر واقعی با میکروفون فیزیکی اندازه‌گیری نشده است.}

برای ارزیابی روی صوت واقعی، رخدادها از مرزهای تفکیک گوینده کشف شدند. از
\N{1,730} رخداد، \N{1,048} رخداد انتخاب و بر پایه نشست تقسیم شد. آستانه
تصمیم تنها روی نشست‌های اعتبارسنجی انتخاب شد و یک‌بار روی آزمون اعمال گردید.

\section{طراحی آزمایش‌های مسیر مستقیم}
\label{sec:direct}

مسیر مستقیم، تطبیق کم‌پارامتر \en{Moshika} با آموزش‌دهنده رسمی است. ورودی،
همان استریوی بخش~\ref{sec:data} است. زیان، ترکیب وزن‌دار زیان متنی و زیان
نشانه‌های صوتی است.

شرط پذیرش زمان اجرا پیش از آموزش قفل شد: آداپتر در سرور جویباری رسمی بار
می‌شود و برای هر یک از ۹ ردیف ازپیش‌تعیین‌شده باید متن، گفتار غیرساکت و قید
خط فارسی برقرار باشد. چک‌پوینت تنها با \N{9/9} واجد شرایط است. این شرط
عمداً سخت‌گیرانه است، چون کمترین زیان اعتبارسنجی در نسخه نخست تقریباً صوت
ساکت و بدون متن می‌داد.

آزمایش‌ها زنجیره فرضیه‌های قابل ابطال‌اند و در هر نسل عمدتاً یک عامل تغییر
کرد (جدول~\ref{tab:moshi-design}). نسخه \en{v6.2} آزمون ظرفیت است نه تعمیم:
آموزش و ارزیابی روی همان \N{32} نمونه، پنجره \N{12} ثانیه، و حذف تصادفی
زمان‌بندی‌شده ورودی متنی از \N{0.25} تا \N{0.75}. اگر زیان حتی روی داده
آموزشی کاهش نیابد مشکل ظرفیت است؛ اگر کاهش یابد اما اجرا شکست بخورد، مشکل
اریبی مواجهه است.

\begin{table}[ht]
\centering
\caption{طراحی شش نسل آزمایشی مسیر مستقیم}
\label{tab:moshi-design}
\small
\begin{tabular}{lL{3.2cm}L{4.6cm}c}
\toprule
\textbf{نسخه} & \textbf{عامل تغییر‌یافته} & \textbf{فرضیه آزمون‌شده} &
\textbf{رتبه} \\
\midrule
\en{v1} & پیکربندی پایه، پنجره \N{20} ثانیه &
مدل با داده کافی گفتار فارسی می‌آموزد & \N{64} \\
\en{v2} & پنجره متن \N{100} ثانیه &
پنجره کوتاه آغاز پاسخ را در بافت نمی‌گنجاند & \N{64} \\
\en{v3} & انجماد درج‌گرها، رتبه بالاتر &
آموزش درج‌گرها بازنمایی پایه را تخریب می‌کند & \N{128} \\
\en{v4} & تنها دو درج‌گر متنی &
تنها فضای متنی فارسی نیاز به تطبیق دارد & \N{128} \\
\en{v5} & ضریب زیان کتاب رمز نخست از \N{100} به \N{10} &
وزن زیاد کتاب رمز معنایی یادگیری آوایی را خفه می‌کند & \N{128} \\
\en{v6.2} & آزمون ظرفیت درون‌نمونه، پنجره \N{12} ثانیه &
مدل ظرفیت یادگیری دارد و مشکل اریبی مواجهه است & \N{64} \\
\bottomrule
\end{tabular}
\end{table}

\clearpage
\section{روش کاری شواهد‌محور}
\label{sec:evidence-first}

پنج قاعده عملیاتی اعمال شد. یک: پروتکل، شاخص ردیف‌ها و آستانه‌ها پیش از
اجرا ثبت می‌شوند. دو: ارزیابی معنایی دو‌مرحله‌ای است؛ تنها پس از عبور کامل
صلاحیت هفت‌ردیفی، آزمون نهایی یک‌بار باز می‌شود. داور
\en{Qwen3.8-27B-FP8} با دمای صفر، ترتیب دوسوکور و دو تکرار است. سه: مصنوعات
نتیجه یک‌بار‌نوشت‌اند. چهار: هر اندازه‌گیری یک کلاس شواهد دارد
(جدول~\ref{tab:evidence-classes}) و تنها کلاس رسمی می‌تواند نیازمندی تعریف
پروژه را اثبات کند. پنج: تجمیع‌کننده شکست‌بسته است؛ نبود مصنوع یا کلاس
پایین‌تر از حد لازم، شرط را ناموفق می‌کند نه نامعلوم.

\begin{table}[ht]
\centering
\caption{کلاس‌های شواهد و توان ادعای هر کلاس}
\label{tab:evidence-classes}
\small
\begin{tabular}{L{3.4cm}L{5.6cm}c}
\toprule
\textbf{کلاس شواهد} & \textbf{شرایط لازم} & \textbf{اثبات نیازمندی} \\
\midrule
رسمی سرتاسری & نشست زنده با رضایت، تأیید پخش کارخواه، میکروفون واقعی و
سخت‌افزار واجد شرایط & مجاز \\
زنده غیر‌رسمی & همان پروتکل، بدون سخت‌افزار یا شرایط مطالعه واجد شرایط &
غیر‌مجاز \\
مؤلفه سمت سرور & زمان‌سنجی تولید یا آشکارساز، بدون اثبات رفتار شنیدنی
کارخواه & غیر‌مجاز \\
تقریب مصنوعی & صوت تولیدشده یا رخداد شبیه‌سازی‌شده & تنها آزمون بازگشتی \\
تقریب خودکار هم‌خانواده & داوری با مدل زبانی از خانواده نامزد & غیر‌مجاز \\
\bottomrule
\end{tabular}

\vspace{0.5em}
{\footnotesize\itshape پیامد این روش آن است که وضعیت شواهد هر شرط مستقل از
آمادگی گزارش برای ارائه و داوری بیان شود. نتیجه نهایی، تفکیک روشن آنچه
ساخته شده از آنچه با شواهد رسمی اثبات شده است.\par}
\end{table}
